\section{Objetivo del Proyecto}

\subsection{Contexto Clínico}
El diagnóstico temprano y preciso de tumores cerebrales constituye
uno de los retos más relevantes de la neuro-oncología
contemporánea. La resonancia magnética (MRI, \textit{Magnetic
Resonance Imaging}) es la modalidad de imagen de referencia para
la detección y caracterización de estas lesiones; sin embargo, su
interpretación depende fuertemente de la experiencia del radiólogo
y está sujeta a variabilidad inter-observador. En este contexto,
los sistemas de apoyo al diagnóstico basados en aprendizaje
profundo han mostrado resultados prometedores como herramientas
de cribado y segunda opinión, capaces de procesar grandes
volúmenes de estudios de manera consistente.

\subsection{Objetivo Técnico}
El presente proyecto tiene como objetivo diseñar, entrenar y
evaluar tres arquitecturas propias de redes neuronales
convolucionales (CNN) para la clasificación automática de
imágenes MRI cerebrales en cuatro categorías:
\textit{glioma\_tumor}, \textit{meningioma\_tumor},
\textit{no\_tumor} y \textit{pituitary\_tumor}. Se utiliza el
dataset público \textit{Brain Tumor Classification MRI} disponible
en Kaggle. De manera explícita, este trabajo \textbf{no} emplea
técnicas de \textit{transfer learning} ni pesos preentrenados
sobre ImageNet u otros corpus; el énfasis pedagógico está en la
construcción de arquitecturas convolucionales desde cero y en la
comprensión del impacto de sus hiperparámetros y bloques
constitutivos.

\subsection{Criterios de Evaluación}
La evaluación de los modelos se realizará mediante un conjunto
de métricas clásicas de clasificación multiclase:
\textit{Accuracy}, \textit{Precision}, \textit{Recall} y
\textit{F1-Score} en su variante \textit{macro-averaged},
complementadas con la matriz de confusión por modelo. Dado el
desbalance entre clases observado durante el análisis exploratorio
(Sección~\ref{sec:eda}), el \textit{F1-macro} se adopta como
métrica principal de comparación, por su robustez frente a clases
minoritarias.
